% Generated draft: requires mathematical and source review.
\documentclass[11pt]{article}
\usepackage[margin=25mm]{geometry}
\usepackage{fontspec}
\setmainfont{Latin Modern Roman}
\usepackage{amsmath,amssymb,mathtools,unicode-math}
\setmathfont{Latin Modern Math}
\usepackage{enumitem,needspace,xurl,hyperref}
\setlength{\parindent}{0pt}
\setlength{\parskip}{5pt}
\setlength{\emergencystretch}{3em}
\newcommand{\N}{\mathbb{N}}
\newcommand{\Z}{\mathbb{Z}}
\newcommand{\Q}{\mathbb{Q}}
\newcommand{\R}{\mathbb{R}}
\newcommand{\C}{\mathbb{C}}
\newcommand{\F}{\mathbb{F}}
\newcommand{\A}{\mathbb{A}}
\newcommand{\PP}{\mathbb{P}}
\newcommand{\E}{\mathbb{E}}
\newcommand{\eps}{\varepsilon}
\newcommand{\dd}{\,\mathrm{d}}
\begin{document}
\section*{UnsolvedMath ID: 800112}
\textbf{Source problem code:} catalogue entry
\subsection{Definitions and mathematical statement}
Let \(S^2\) be the oriented two-sphere, and write \(\mu\) for the area measure intended by the source. The extracted record does not specify the area form or measure. Define
\[
G=\{h:S^2\to S^2\mid h\text{ is an orientation-preserving homeomorphism and }\mu(h(A))=\mu(A)\text{ for every Borel set }A\subseteq S^2\}.
\]
For \(f,g\in G\), define the commutator by \([f,g]=fgf^{-1}g^{-1}\), and let
\[
D=[G,G]=\left\{\prod_{i=1}^{n}[f_i,g_i]\;\middle|\;n\in\Z_{\geq 0},\ f_i,g_i\in G\right\},
\]
where the empty product is the identity. Thus \(D\) is the algebraic commutator subgroup; no topological closure is taken. The abelianization is the abstract quotient group \(G^{\mathrm{ab}}=G/D\).

The supplied source reports that \(D\subsetneq G\), equivalently that \(G\) is not perfect. This is a reported background result, not an independently verified conclusion of this draft.

\textbf{Question 53.} Determine the commutator subgroup \([G,G]\) and the abelianization \(G/[G,G]\). The source does not specify a required form of description for either group, and supplies no topology for the quotient. Its intended area structure should be confirmed before a definitive formulation is accepted.
\subsection{Short English statement}
For the group of area- and orientation-preserving homeomorphisms of the two-sphere, describe the subgroup generated by commutators and the resulting abelian quotient. The source reports that this subgroup is proper and labels both questions open. The intended area measure and the desired form of description remain unspecified.
\subsection{Sources}
\begin{itemize}
\item Local UnsolvedMath registry, permanent ID \(800112\), exact source problem code \texttt{catalogue entry}; title: \emph{Foliations and Diffeomorphism Groups --- Question 53}. The supplied record attributes the problem to Nariman and Yazdi (2025) and labels it \texttt{open}.
\item Nariman, Yazdi, \emph{Problem List on Foliations and Diffeomorphism Groups} (2025), Question 53. Supplied source URL: \url{https://arxiv.org/abs/2503.18274}. The record gives access date \texttt{2026-10-16} and extraction method \texttt{arxiv\_latex}; these are supplied provenance details, not evidence of an independent source visit for this draft.
\item For the reported nonperfectness result, the extracted statement cites the reference labels \texttt{cristofaro2023pfh}, \texttt{cristofaro2024proof}, and \texttt{cristifaro-gardiner2022quantitative}. Complete bibliographic entries and links for these labels are not supplied, so their identities and support for the reported result remain unverified.
\end{itemize}
\end{document}
